% Generated draft: requires mathematical and source review.
\documentclass[11pt]{article}
\usepackage[margin=25mm]{geometry}
\usepackage{fontspec}
\setmainfont{Latin Modern Roman}
\usepackage{amsmath,amssymb,mathtools,unicode-math}
\setmathfont{Latin Modern Math}
\usepackage{enumitem,needspace,xurl,hyperref}
\setlength{\parindent}{0pt}
\setlength{\parskip}{5pt}
\setlength{\emergencystretch}{3em}
\newcommand{\N}{\mathbb{N}}
\newcommand{\Z}{\mathbb{Z}}
\newcommand{\Q}{\mathbb{Q}}
\newcommand{\R}{\mathbb{R}}
\newcommand{\C}{\mathbb{C}}
\newcommand{\F}{\mathbb{F}}
\newcommand{\A}{\mathbb{A}}
\newcommand{\PP}{\mathbb{P}}
\newcommand{\E}{\mathbb{E}}
\newcommand{\eps}{\varepsilon}
\newcommand{\dd}{\,\mathrm{d}}
\begin{document}
\section*{UnsolvedMath ID: 20001894}
\textbf{Source problem code:} catalogue entry
\subsection{Definitions and mathematical statement}
The extracted source formula is fragmented. Its coherent working reconstruction is
\[
 N(T)=\bigl(1+O(e^{-\varepsilon T})\bigr)
       \operatorname{Li}(e^{hT}),\qquad
 \operatorname{Li}(x)=\int_2^x\frac{du}{\log u},
\]
for some \(\varepsilon>0\), where \(h>0\) is topological entropy and \(N(T)\) counts primitive oriented periodic geodesic-flow orbits of period at most \(T\). The source omits the preceding dynamical/metric assumptions and the counting convention. The aid's transfer statement concerns a compact hyperbolic/Anosov geodesic setting with finite bounded-period counts and a positive shortest period, not every metric on every surface. The error constant and exponent need not be uniform over metrics.

For a reversible Riemannian geodesic flow, each unoriented geometric prime geodesic gives two oriented flow orbits. Hence the exact convention change is
\[
 N_{or}(T)=2N_{un}(T),
\]
and the unoriented main coefficient is \(1/2\), not one. This factor is not automatically applicable to nonreversible Finsler flows. Counting proper iterates is a separate change: with \(I(T)=\sum_{k\ge1}N(T/k)\), the aid reports
\[
 I(T)-N(T)=O(e^{hT/2}/T),
\]
generally sharp because double iterates already contribute \(N(T/2)\) of that order under the usual prime asymptotic. Thus ignoring repeats can impose an \(h/2\) relative-error barrier even when primitive counting has a better estimate.

A precise weighted normalization for the reported Chebyshev transfer is
\[
 \Psi(T)=\sum_{p\text{ primitive}}\ \sum_{k\ell(p)\le T}\ell(p).
\]
Assume for some \(\delta>0\),
\[
 \Psi(T)=e^{hT}/h+O(e^{(h-\delta)T}).
\]
The aid reports that removing proper iterates and partial summation yield the reconstructed AIM formula with
\[
 \varepsilon=\min\{\delta,h/2\}.
\]
The weight is the primitive length \(\ell(p)\), not the iterated length \(k\ell(p)\). The estimate is a conditional input; the convention-transfer lemma does not prove the spectral gap required to obtain it.

The source says earlier results cover the intended metrics with \(C^1\) stable/unstable bundles, mentioning surfaces and quarter-pinched metrics. Exact curvature, pinching strictness, regularity, and mixing hypotheses are not supplied, so this is retained as a source assessment rather than broadened to all surface metrics. The aid also gives a dated literature warning: the GLP erratum invalidates the original broader contact-bunching spectral-gap proof and restores it only under a stronger homogeneity hypothesis. It says arXiv:2501.02537v10 claims the full contact-Anosov case but is a substantially corrected preprint. No unconditional full-case solution is inferred from that claim, and no exact erratum hypotheses or complete bibliography are supplied. The partial result here is the conditional counting/convention conversion, not a resolution of the whole conjecture.
\subsection{Short English statement}
The aid reconstructs the fragmented conjecture as a primitive oriented prime-orbit count with relative exponential error. Its transfer lemma fixes the orientation factor, proper-iterate contribution, and Chebyshev normalization, while the broad spectral-gap input remains explicitly qualified by the erratum and preprint caveats.
\subsection{Sources}
\begin{itemize}
\item Local UnsolvedMath registry, permanent ID \(20001894\), source code \texttt{catalogue entry}; \texttt{aim-geometry-notes.json}, \texttt{notes[231]}, zero-based index 231, tag \texttt{conjecture}.
\item AIM workshop ``Geodesics,'' recorded source item 5.5, extracted Conjecture 5.5.1: \url{https://aimath.org/pastworkshops/geodesicsproblems.pdf}. Source type PDF; statement status \texttt{exact} for the supplied extraction, though its formula is fragmented. Only that supplied text was used; no PDF was accessed or compiled.
\item Earlier surface/regularity/pinching results and the GLP erratum are mentioned without authoritative full citations or exact hypotheses. The aid cites the substantially corrected preprint arXiv:2501.02537v10: \url{https://arxiv.org/abs/2501.02537v10}, as claiming the full contact-Anosov case. That claim and the erratum assessment have not been independently verified. Original remarks and literature fields are empty.
\item Research attempt 1: \texttt{partial\_result}, status at run \texttt{partially\_solved}, rigor label \texttt{proved}, novelty confidence low; registry status \texttt{partially\_solved}, classification \texttt{PARTIAL-PROGRESS}. Created and updated 14 August 2026. Machine-generated aid, not peer reviewed; independently verify before citation.
\item Ulam AI / UnsolvedMath Contributors; CC BY 4.0.
\end{itemize}
\end{document}
